\chapter{Conclusion}\label{chapter:conclusion}

This dissertation is centered around giving a comprehensive accounting of the 
current approaches to mediation analysis with multiple mediators. We 
discussed the definitions, identification assumptions, and interpretation of natural and 
interventional effects, the two most popular approaches within the mediation analysis 
literature. We first reviewed why it has generally been established that identification
of natural effects is infeasible for mediation analysis with multiple mediators. Having 
motivated the impetus to search for alternative mediational effect definitions, we then 
built up the theory for interventional effects for one and two mediators, with a 
particular focus on defining hypothetical experiments that could correspond to the 
interventions implied by these effects. We argued that though these effects are 
theoretically easier to identify, they cannot be interpreted in the same way as 
natural effects because they do not estimate the strength of causal mechanisms. 
Additionally, it is difficult to find applications where the value of these 
effects --- estimating the causal effects of particular joint stochastic 
interventions --- can be clearly seen. 

Given the problems with both natural and interventional effects in multiple mediator 
settings, we then propose VanderWeele-Robinson interventional effects as a 
promising class of mediational effects with clear real-world applications. We 
formalize and extend these effects to two mediator problems. We then apply 
these effects to a real-world dataset using a flexible Monte Carlo estimation 
strategy and discuss how the practical implications of these effect estimates can 
be thought about. 

A key argument made throughout this dissertation is that there is little conceptual 
overlap between the interpretation of natural and interventional effects. This means 
that important features of natural effect definitions --- like the ability to decompose
total causal effect into path-specific causal mechanisms --- do not hold the same 
importance for the interpretation of interventional effects. Despite this, the 
definition of interventional effects for multiple mediator problems in particular 
seems to have been influenced by attempts to mimic the features of natural effect 
definitions. Given the recent work of~\cite{miles2023}, which the core arguments of 
this dissertation owe a large debt to, we believe that interventional
effects cannot be seen as a ``next-best'' substitute for natural effects and that 
future work on the theory behind interventional effects should be motivated by 
the unique prescriptive interpretation of these effects.

The are a number of potentially interesting directions for future work related 
to the topics of this dissertation. On the development of interventional effects for 
multiple mediators, it would be interesting to further explore what kinds of unique 
real-world applications such stochastic joint interventions can have. This 
dissertation shows that VanderWeele-Robinson style interventional effects already 
offer one such promising variation. While these effects have recently been popularized 
in the health disparities literature~\cite{jackson2021}, much interested work 
remains to be done. Given that the literature on these effects has so for 
proposed parametric estimation strategies, it would be interesting to see to what 
extent the semiparametric estimation theory for interventional effects carries 
over to the efficient estimation of VanderWeele-Robinson effects. Additionally, there
are potentially interesting connections between the population level intervention 
framework of~\cite{hubbardvanderlaan2008} and VanderWeele-Robinson effects. Finally, 
Sections 6 of this dissertation considered an application of VanderWeele-Robinson 
effects under the simplifying assumption of no covariates. It would be interesting to 
combine the insights of Section~\ref{subsection:vr_effects_2m} 
and~\cite{jacksonvanderweele2018} to define these effects for multiple mediator situations
that properly account for pre-treatment covariates. 
